(1) 概率的构象熵给位移相关自由能 $\Delta F=-Nk_BT\ln[P(x)/P(0)]=3Nk_BT x^2/(2nb^2)$。此式只描述Gaussian构象统计，忽略高伸长时有限轮廓长度非线性。

(2) 外部为保持位移所需的力 $f=\partial_x\Delta F=3Nk_BTx/(nb^2)$。构象内能变化 $\Delta U=\Delta F-T\partial_T\Delta F=0$，而熵差 $\Delta S=-\partial_T\Delta F=-3Nk_B x^2/(2nb^2)$。所以弹性力本质是熵恢复力，$f\propto T$。

(3) 由 $f=3Nk_BTx/(nb^2)$ 得 $(\partial f/\partial T)_x=3Nk_Bx/(nb^2)$，而 $\partial_x(-\Delta S)=3Nk_Bx/(nb^2)$，两边严格一致。恒温拉伸使系统构型熵下降，环境必须接收热量 $Q_{\rm chain}=T\Delta S<0$（以入链热为正），外加机械功转为释放的热。

(4) 完整系统熵可写 $S(T,x)=C\ln(T/T_*)-3Nk_Bx^2/(2nb^2)+S_*$。可逆绝热 $S_f=S_i$ 意味着 $C\ln(T_f/T_i)=[3Nk_B/(2nb^2)](x_f^2-x_i^2)$。故 $T_f=T_i\exp[3Nk_B(x_f^2-x_i^2)/(2nb^2C)]$，若拉伸使 $x_f^2>x_i^2$，则 $T_f>T_i$。
